\section{Graphs with Cut-Value at Most Two}
\label{sec:cut-two}

We now turn to another structural property that leads to the same simple formula.
Specifically, we consider graphs in which the minimum cut-value of every pair of vertices is at most two.
This condition restricts the connectivity structure of the graph in a very particular way.

Such graphs have the structure of a \emph{cactus graph} (cf.~\cite{Mehlhorn2017Certifying}): they consist of cycles and single edges, any two of which share at most one vertex, an articulation point.
Indeed, if two distinct cycles shared an edge, their union would contain two vertices joined by three edge-disjoint paths; separating these two vertices would require at least three edges, contradicting the assumption $c(x,y) \leq 2$.
This structure leads to the following result.

\begin{theorem}
\label{thm:cut-two}
Let $G = (V, E)$ be a connected graph such that $c(x, y) \leq 2$ for every pair of distinct vertices $x, y \in V$.
Then, for any two distinct vertices $u,v \in V$, the value of the minimum cut-path is
\[
\cp(u, v) = c(u, v) + d(u, v) - 1.
\]
\end{theorem}

\begin{proof}
We distinguish two cases based on the values of $d(u, v)$ and $c(u, v)$.

\begin{enumerate}
    \item \textbf{Case} $d(u, v) = 1$ \textbf{or} $c(u, v) = 1$:

    The formula holds by Lemma~\ref{lem:basic-bounds}.

    \item \textbf{Case} $d(u, v) \geq 2$ \textbf{and} $c(u, v) = 2$:

    By Lemma~\ref{lem:basic-bounds}, we have
    \[
    d(u, v) \leq \cp(u, v) \leq d(u, v) + 1.
    \]
    Assume for contradiction that $\cp(u, v) = d(u, v)$.
    Then a minimum cut-path is a shortest $u$--$v$ path $P$ that contains a $u$--$v$ cut, so the graph $G \setminus P$ contains no $u$--$v$ path.

    Since $c(u, v) = 2$, no single edge separates $u$ from $v$; hence no edge of $P$ is a bridge, and every edge of $P$ lies on a cycle of the cactus graph $G$.
    The path $P$ therefore passes through a sequence of cycles $C_1, C_2, \dots, C_t$, and in each cycle $C_i$ it follows one of the two arcs between the vertex where it enters $C_i$ and the vertex where it leaves $C_i$.
    The other arc of each cycle $C_i$ is edge-disjoint from $P$, and the concatenation of these arcs connects $u$ and $v$ in $G \setminus P$.
    This is a contradiction.
    Therefore,
    \[
    \cp(u, v) = d(u, v) + 1 = c(u, v) + d(u, v) - 1. \qedhere
    \]
\end{enumerate}
\end{proof}

By Remark~\ref{rem:algorithm}, a minimum cut-path in such a graph is again obtained as the union of a minimum cut and a shortest path.

